\section{Experimental Setup}\label{sec:exp}
\subsection{Data and conditions}
We vary \emph{occlusion}, using 3, 4, 5 or 6 foreground branches, and \emph{viewpoint spacing}, using 20, 40 or 60~mm. Each branch count has one rail capture of 28 views at 20~mm steps (27 at six branches). The 40 and 60~mm conditions keep every second or third view around the same reference, so that spacing changes while the scene and reference stay fixed. Images are processed at 0.8 of full resolution ($2419\times3226$~px, width by height; the phone was held in portrait). Fruit and foliage masks come from SAM~3 and were stored with each capture. All conditions share one physical arrangement of red grapes. The study therefore characterises these conditions, not a population of scenes.

\textbf{Ground truth.} After the foliage was removed, a leaf-free photograph of the same fruit was taken from a different position. Its fruit mask comes from SAM~3, and we hand-labelled the centres of its 25 grapes. For each capture, the photograph is registered to the reference view with a similarity transform fitted to SIFT features near the fruit. We inspected an overlay for every branch count. Agreement is good overall, but some clusters at five and six branches show visible offset, because the fruit lies at different depths.

\subsection{Metrics}
\textbf{Pixel-wise.} With registered truth fruit mask $G$ and recovered mask $\hat G$, we report recall $R=|G\cap\hat G|/|G|$, precision $P=|G\cap\hat G|/|\hat G|$ and IoU. The single-view baseline applies the same formulas to the reference view's own fruit mask. Because $G$ also comes from SAM~3, this evaluation shares a model with the method.

\textbf{Count-based.} Pixel overlap cannot tell whether recovered fruit is still separable into berries. We therefore run the GrapeSAM point model~\cite{tong2026vivid}, with its published weights and post-processing (confidence 0.05), on three fruit-only images: the occluded reference view, the reconstruction and the registered leaf-free photograph. None of its settings were tuned on our data. Detections are matched one-to-one to the 25 labels, closest first, within 24~px, which is about one berry radius. We report count F1, the number of detections $n$, and average precision (AP) over a sweep of confidence thresholds. We call the leaf-free result a \emph{comparison point} rather than a ceiling, because GrapeSAM overcounts it (29--40 detections for 25 grapes).

\textbf{Uncertainty.} With one scene, we cannot measure variation between scenes, so we estimate three narrower sources: which grapes happen to be present (2,000 bootstrap resamples of the 25 grapes, paired between the single view and the reconstruction); the choice of reference (the 20~mm runs repeated with four other references, offset by $\pm4$ and $\pm8$ views); and truth registration error (the truth transform perturbed by shifts of 1--5~px in eight directions, rotations of $\pm0.5^\circ$ and $\pm1^\circ$ and scale changes of $\pm1\%$, then all metrics recomputed).

\subsection{Baselines and ablations}
All baselines use the same views, masks and evaluation. \emph{Mask voting} takes the union, or the majority, of the aligned SAM~3 fruit masks, coloured by the median observed fruit colour. \emph{Synthetic aperture} (SA)~\cite{vaish2006} takes the per-pixel mean, median or foliage-masked mean of the aligned stack, then segments the fused image with SAM~3. We test SA under two alignments. \emph{Depth shift}, our prototype (Sec.~\ref{sec:evolution}), applies one horizontal shift per rail step, chosen by maximising focus over the reference fruit region. \emph{Our alignment} uses the jointly solved poses, which isolates the effect of the fusion rule. Component ablations use five branches at 20~mm, where the gain is largest. Fusion and mask variants reuse the full method's poses. Alignment variants re-solve the poses and pass them through the same validation.

To separate spacing from the number of views, we vary the number of views centred on the reference at 20~mm, from 3 to 28. We also fix nine views at 20, 40 and 60~mm, and fix the span at 480~mm, which takes 25, 13 and 9 views at those spacings.

\textbf{Implementation.} The code is in Python with OpenCV~\cite{opencv}, NumPy~\cite{numpy}, Pillow and scikit-image; SAM~3~\cite{sam3code} and GrapeSAM~\cite{tong2026vivid} run under PyTorch~\cite{pytorch} on the GPU. The pipeline, the mask scripts and the scripts that produce the tables and figures are in our code repository~\cite{teamrepo}. As a check of the code, a repeat run on one capture reproduced the same reference view, the same accept or reject decision for 101 of the 102 tested pairs and the same 13,395 observed pixels.
